% !TEX root = ../main.tex
\title{Problem Solving - Boulos 3 ed.}
\author{Túlio Gomes}
\date{may 2026}

\maketitle

\section{Introduction}
This is a document that contains all solutions i'll make on the Boulos 3 ed. book. The main goal is to practice and improve my problem-solving skills in analytical geometry, and to have a reference for future studies.

\section{Chapter 22 - Ellipse, Hyperbola and Parabola}
\noindent
% Definition 22-1 Section
\newcommand{\ellipseDefinition}{d(X,F_1) + d(X,F_2) = 2a}
\noindent
\begin{tcolorbox}[colback=gray!5, colframe=gray!60, arc=2mm, title=\textbf{22-1 \quad Defini\c{c}\~{a}o}, coltitle=white, fonttitle=\bfseries]
Sejam $F_1$ e $F_2$ pontos distintos, $2c$ sua dist\^{a}ncia e $a$ um n\'{u}mero real tal que $a > c$. O lugar geom\'{e}trico $\mathbf{E}$ dos pontos $X$ tais que $d(X,F_1) + d(X,F_2) = 2a$ chama-se \textbf{elipse}. Cada um dos pontos $F_1$ e $F_2$ \'{e} chamado \textbf{foco} da elipse, o segmento $F_1F_2$ \'{e} chamado \textbf{segmento focal}, seu ponto m\'{e}dio, \textbf{centro} da elipse, e $2c$, \textbf{dist\^{a}ncia focal}. A reta $F_1F_2$ chama-se \textbf{reta focal}, e qualquer segmento cujas extremidades (distintas) pertencem a $\mathbf{E}$ chama-se \textbf{corda} da elipse.
\end{tcolorbox}

\vspace{2em}

% Exercise 22-2 Section
\textbf{22-2} \quad Sejam $\mathbf{E}$ e $a$ como na Definição 22-1. Prove que:
\begin{enumerate}[label=(\alph*), leftmargin=2.5em, itemsep=0.8ex, parsep=0ex, topsep=0.5ex]
    \item os focos e o centro de $\mathbf{E}$ não pertencem a $\mathbf{E}$;
    \item nenhum ponto do segmento focal de $\mathbf{E}$ pertence a $\mathbf{E}$;
    \item existem precisamente dois pontos $A_1$ e $A_2$ da reta focal de $\mathbf{E}$ que pertencem a $\mathbf{E}$;
    \item[\llap{$\rightarrow$\hspace{0.4em}}(d)] se $PQ$ é uma corda qualquer de $\mathbf{E}$, então $d(P,Q) \le 2a$;
    \item[\llap{$\rightarrow$\hspace{0.4em}}(e)] a única corda de $\mathbf{E}$ de comprimento $2a$ é $A_1A_2$.
\end{enumerate}

\textbf{Solução de (a)}

Colocando, convenientmente, o centro em (0,0) e o eixo das abscissas sobre $F_1F_2$, temos: $O = (0, 0)$; $F_1 = (-c, 0)$; $F_2 = (c, 0)$.

Portanto, segue a equação da elipse, a partir de $X=(x, y)$, com os parâmetros atualizados:
\newcommand{\distanceDefinitions}{
\begin{aligned}
d(X,F_1) &= \sqrt{(-c - x)^2 + (-y)^2} \\
d(X,F_2) &= \sqrt{(c - x)^2 + (-y)^2}
\end{aligned}
}
\begin{equation}
\distanceDefinitions
\label{eq:distance-definitions}
\end{equation}
\begin{equation}
    \begin{split}
        \sqrt{(-c-x)^2 + (-y)^2} + \sqrt{(c-x)^2 + (-y)^2} = 2a \\
        \frac{ \sqrt{(-c-x)^2 + (-y)^2} + \sqrt{(c-x)^2 + (-y)^2}}{2} = a
    \end{split}
\end{equation}
Se a > c, então:
\begin{equation}
    \frac{ \sqrt{(-c-x)^2 + (-y)^2} + \sqrt{(c-x)^2 + (-y)^2}}{2} > c
\end{equation}
De acordo com a hipótese, os focos e o centro não pertencem a $\mathbf{E}$, ou seja, não satisfazem a equação que define $\mathbf{E}$. Logo, $F_1=(-c,0)$ não satisfaz a equação:
\begin{equation}
    \begin{split}
        \frac{\sqrt{(-c-(-c))^2 + (0)^2} + \sqrt{(c-(-c))^2 + (0)^2}}{2} > c \\
        \frac{\sqrt{(c+c)^2}}{2} > c \\
        c > c
    \end{split}
\end{equation}
O mesmo ocorre para $F_2=(c,0)$:
\begin{equation}
    \begin{split}
        \frac{\sqrt{(-c-c)^2 + (0)^2} + \sqrt{(c-(c))^2 + (0)^2}}{2} > c \\
        \frac{\sqrt{4c^2}}{2} > c \\
        c > c
    \end{split}
\end{equation}
Assim como para O:
\begin{equation}
    \begin{split}
        \frac{\sqrt{(-c-0)^2 + (0)^2} + \sqrt{(c-0)^2 + (0)^2}}{2} > c \\
        \frac{\sqrt{c^2} + \sqrt{c^2}}{2} > c \\
        c > c
    \end{split}
\end{equation}